\subsubsection{规划器}

规划器（全局规划器）是ROS中的核心模块之一，它负责规划机器人的全局路径，即从起始点到目标点的路径。
他们通过全局代价地图global\_costmap，根据当前机器人的位置和目标位置，计算出从当前位置到目标位置的(时间或路程)最优路径。
路径规划算法有很多种，还有他们的优化版本。
这里，我们介绍两个最经典好用的算法：$Dijkstra$算法和$A*$算法
\paragraph{$Dijkstra$算法}
$Dijkstra$算法是一种最简单的广度优先的路径规划算法，它通过计算从起始点到其他所有点的最短路径，直到覆盖到终点。

$https://www.bilibili.com/video/BV19b4y1d7Hz/$


\subsubsection{控制器}
如果只有规划器会怎么样呢？现在，登录你舍友的LOL账号，开一把排位晋级赛，在小地图上点击敌方基地。
好的，一条全局导航路径生成出来了；好的，你的英雄沿着规划好的路径动起来了；好的，你被对面拿下一血并被队友问候全家。
这是为什么？因为敌方是动态的，他的位置不是事先知道的，事先知道的只有地图，而全局规划器只考虑了地图，未考虑动态障碍物。
局部规划器就是在你的视野里出现动态障碍物的时候，规划出一条道路绕开他，然后继续跟着全局路径走。

回到缺德地图，你这位车手就是局部规划器，缺德地图（全局规划器）只给大致的路线。
如何避让冲出来的小学生，掉头，超车，等红灯等都是车手（局部规划器）决定的。

常用的局部规划器有DWA,TEB等，我们使用的是TEB算法。